

%__________________________________________________________________________________
% \logametrique{BG}{MD}{10}{0}{0.8}{black}{0.08}
% BG  HD  x10    retrait(cm) couleur  resolution (cm)
% __________________________________________________________________________________
\newcommand{\logametrique}[7]{
  \coordinate (Z) at (0,0) {};
  \coordinate (U) at (Z|-#1) {}; %disons le 1 de l'échelle
  \coordinate (D) at ($(U)+(#3,#4)$) {}; %disons le 1 de l'échelle
  \node[] at ($(U)$){1};
  \node[] at ($(D)$){10};
  
  %% On calcule la résolution la plus fine
  % \pgfmathsetmacro{\dep}{int(-(log10(10^(#7/#3)-1))+1)}
  \pgfmathsetmacro{\dep}{int(-log10( 10^(#7/#3) -1 ))}
  \pgfmathsetmacro{\deppun}{\dep + 1}
  % \pgfmathsetmacro{\dep}{2}
  \pgfmathsetmacro{\sep}{#5/(\dep+3))}
  
  % \foreach \d  in {2,...,2} {
  % \coordinate (Q) at ($(U)!{0.1*\dep}!(D)$) {};
  % \pgfmathsetmacro{\kmax}{ int( (1./(1.-10^(-#7/#3))-10^(\dep))/10. )*10}
  \pgfmathsetmacro{\kreso}{    int( (1 + 1 / (10^(#7/#3)-1) ) / 10 ) * 10}
  \pgfmathsetmacro{\kdeuxreso}{int( (1 + 2 / (10^(#7/#3)-1) ) / 20 ) * 20}
  \pgfmathsetmacro{\kcinqreso}{int( (1 + 5 / (10^(#7/#3)-1) ) / 50 ) * 50}
  \pgfmathsetmacro{\kmax}{10^(\dep)*9-1}
  \pgfmathsetmacro{\kun}  {(\kreso    >\kmax  )?\kmax  :( \kreso    <1 ? 1 : \kreso   )}
  \pgfmathsetmacro{\kdeux}{(\kdeuxreso>\kmax-1)?\kmax-1:( \kdeuxreso<1 ? 1 : \kdeuxreso )}
  \pgfmathsetmacro{\kcinq}{(\kcinqreso>\kmax-4)?\kmax-4:( \kcinqreso<1 ? 1 : \kcinqreso )}
  \pgfmathsetmacro{\kunp}{\kun+2}
  %\node[] at ($(0,1)$){DEP = \dep};
  %\node[] at ($(0,2)$){kmax = \kmax};
  %\node[] at ($(0,3)$){kreso = \kreso};
  %\node[] at ($(0,4)$){k2 = \kdeuxreso};
  %\node[] at ($(0,5)$){k5 = \kcinqreso};
  
  \foreach \k in {0,...,\kmax}{
    \pgfmathparse{(\k<=\kreso) && (modulo(\k,2)!=0)  && (modulo(\k,5)!=0) ? 1 : 0}
    \ifthenelse{\pgfmathresult = 1}{
      \coordinate (Q) at ($(U)!{logpos(\k,\dep)}!(D)$) {};
      % \node[] at ($(Q)+(0,\dep*\sep+3*\sep-0.2)$){\k}; 
      \draw[shift=(Q), xstep=#3cm, ystep=#4cm, line width=1pt, #6!100] ($(#1)+(0,\dep*\sep+3*\sep)$) grid (#2);
      
    }{
      \pgfmathparse{(\k<\kdeuxreso) && (modulo(\k,2)==0 && (modulo(\k,5)!=0) ? 1 : 0}
      \ifthenelse{\pgfmathresult = 1}{
        \coordinate (Q) at ($(U)!{logpos(\k,\dep)}!(D)$) {};
        % \node[] at ($(Q)+(0,\dep*\sep+2*\sep-0.2)$){\k}; 
        \draw[shift=(Q), xstep=#3cm, ystep=#4cm, line width=1pt, #6!100] ($(#1)+(0,\dep*\sep+3*\sep)$) grid (#2);
        
      }{
        \pgfmathparse{(\k<\kcinqreso) && ((\k>\kdeuxreso)||(\k<\kreso))  && (modulo(\k,2)!=0 && (modulo(\k,5)==0)  ? 1 : 0}
        \ifthenelse{\pgfmathresult = 1}{
          \coordinate (Q) at ($(U)!{logpos(\k,\dep)}!(D)$) {};
          % \node[] at ($(Q)+(0,\dep*\sep+\sep-0.2)$){\k}; 
          \draw[shift=(Q), xstep=#3cm, ystep=#4cm, line width=1pt, #6!100] ($(#1)+(0,\dep*\sep+\sep)$) grid (#2);
          
        }{
          \pgfmathparse{(modulo(\k,10)==0  ? 1 : 0}
          \ifthenelse{\pgfmathresult = 1}
          {
            \coordinate (Q) at ($(U)!{logpos(\k,\dep)}!(D)$) {};
            \pgfmathsetmacro{\dep}{int(-(log10(10^(#7/#3)-1))+1)}
            
            
            % \node[] at ($(Q)+(0,\dep*\sep-0.2)$){\k}; 
            \ifthenelse{\k = 0}
            {
              \draw[shift=(Q), xstep=#3cm, ystep=#4cm, line width=2pt, #6!100] ($(#1)+(0,\sep)$) grid (#2);
              % \node[] at (#2-|Q){\k};       
            }
            {
              \pgfmathparse{modulo(\k, 10^(\dep-1)) == 0 ? 1 : 0}
              \ifthenelse{\pgfmathresult = 1}
              {
                \pgfmathsetmacro{\puisdix}{\dep-int(log10(\k)+1)}
                \draw[shift=(Q), xstep=#3cm, ystep=#4cm, line width=2pt, #6!100] ($(#1)+(0,\sep)$) grid (#2);
                \pgfmathparse{int(int(\k/10^(\dep-1))+1)}
                
                \node[] at (Q){\pgfmathresult};
              }
              {
                \pgfmathsetmacro{\puisdix}{\dep-int(log10(\k)+1)}
                \draw[shift=(Q), xstep=#3cm, ystep=#4cm, line width=1pt, #6!100] ($(#1)+(0,\puisdix*\sep)$) grid (#2);
                % \node[] at (#2-|Q){\k};
                
              }
              
            }
          }
          {
            
            
          }%Fin dep+
          
          
        }%Fin cinq
        
      }%Fin deux
    }%Fin uns
  }
  % \foreach \k in {\kun,\kunp,...,\kdeux}{
  % \coordinate (Q) at ($(U)!{log10(\k/(10^\d))}!(D)$) {};
  % \draw[shift=(Q), xstep=#3cm, ystep=#4cm, line width=3-\d, #6!100] ($(#1)+(0,\d*3*\sep-\sep)$) grid (#2);
  % }
  % }
  
  
  %%   2,3,..,9
  %   \foreach \i in {2,...,9} {
  %   \coordinate (P) at ($(U)!{log10(\i)}!(D)$) {};
  %   \draw[shift=(P), xstep=#3cm, ystep=#4cm, line width=0.6mm, #6!75] ($(#1)+(0,0.5*#5)$) grid (#2);
  % }
  
  %%   1, 10, 100
  %   \draw[xstep=#3cm, ystep=#4cm, line width=0.8mm, #6!100] (#1) grid (#2);
  
}



% \logametrique{BG}{MD}{10}{0}{0.8}{black} 
% BG  HD  x10    retrait couleur 
\newcommand{\loglogametrique}[6]{
  \coordinate (S) at ($(#3,#4)$) {}; %disons le 1 de l'échelle
  
  %% 1.1, 1.2, ...  ou 1.2, 1.4, ...  ou 1.5 selon def 
  \foreach \i in {1,...,9} {
    
    %% IF 1.
    \pgfmathparse{(log10(\i+1)-log10(\i))*#3>1?1:0}
    \ifthenelse{\pgfmathresult = 1}{
      %% Then 1.1, 1.2, ... passe
      
      %% => Trace les 1.1 , 1.2, ...
      \foreach \j in {1,...,9} {
        
        %% IF 1.1.1
        
        \pgfmathparse{3>1?1:0}
        % \pgfmathparse{(log10(\i+\j/10)-log10(\i+(\j+1)/10))*#3>1?1:0}
        \ifthenelse{\pgfmathresult = 1}{
          %% Then 1.1.1, 1.1.2, ... passe
          \foreach \k in {1,...,9} {
            \coordinate (Q) at ($(0,0)!{log10(\i+(\j/10)+(\k/100))}!(S)$) {};
            \draw[shift=(Q), xstep=#3cm,  ystep=#4cm, line width=0.2mm, #6!25]  ($(#1)+(0,#5)$) grid (#2);
            
          }%% FIN for k
          
          %% => le 1.1.5
          \coordinate (P) at ($(0,0)!{log10(\i+(\j/10)+.05)}!(S)$) {};
          \draw[shift=(P), xstep=#3cm,  ystep=#4cm, line width=0.4mm, #6!50] ($(#1)+(0,0.*#5)$) grid (#2);
          
        }{
          %% Else passe pas
        } %% FIN IF 1.1.1
        
        
        \coordinate (Q) at ($(0,0)!{log10(\i+(\j/10))}!(S)$) {};
        \draw[shift=(Q), xstep=#3cm,  ystep=#4cm, line width=0.2mm, #6!25]  ($(#1)+(0,1.5*#5)$) grid (#2);
        
        
      } %% FIN FOR J
      %% => le 1.5
      \coordinate (P) at ($(0,0)!{log10(\i+0.5)}!(S)$) {};
      \draw[shift=(P), xstep=#3cm,  ystep=#4cm, line width=0.4mm, #6!50] ($(#1)+(0,#5)$) grid (#2);
      
    }{
      %% ELSE 1.1
      %% IF 1.2
      \pgfmathparse{(log10(\i+1)-log10(\i))*#3>.5?1:0}    
      \ifthenelse{\pgfmathresult = 1}{
        %% THEN 1.2, 1.4,... passe
        \foreach \j in {2,4,6,8} {
          \coordinate (Q) at ($(0,0)!{log10(\i+(\j/10))}!(S)$) {};
          \draw[shift=(Q), xstep=#3cm,  ystep=#4cm, line width=0.3mm, #6!40]  ($(#1)+(0,1.5*#5)$) grid (#2);
        }
      }{
        %% ELSE que 1.5 passe
        \coordinate (P) at ($(0,0)!{log10(\i+0.5)}!(S)$) {};
        \draw[shift=(P), xstep=#3cm,  ystep=#4cm, line width=0.4mm, #6!50] ($(#1)+(0,#5)$) grid (#2);
      } 
    }%%FIN IF 1.
    
    %% 2,3,..,9
    \foreach \i in {2,...,9} {
      \coordinate (P) at ($(0,0)!{log10(\i)}!(S)$) {};
      \draw[shift=(P), xstep=#3cm, ystep=#4cm, line width=0.6mm, #6!75] ($(#1)+(0,0.5*#5)$) grid (#2);
    }
    
    %% 1, 10, 100
    \draw[xstep=#3cm, ystep=#4cm, line width=0.8mm, #6!100] (#1) grid (#2);
    
  }
  
  
  % \foreach \i in {3,...,6} {
  % \foreach \j in {2,4,...,8} {
  % \coordinate (Q) at ($(0,0)!{log10(\i+(\j/10))}!(S)$) {};
  % \draw[shift=(Q), xstep=#3cm, ystep=#4cm, line width=0.2mm, #6!25] ($(#1)+(0,1.5*#5)$) grid (#2);
  % }
  % }
  
  %   \foreach \i in {1,2,7,8,9} {
  %   \coordinate (P) at ($(0,0)!{log10(\i+0.5)}!(S)$) {};
  %   \draw[shift=(P), xstep=#3cm,  ystep=#4cm, line width=0.4mm, #6!50] ($(#1)+(0,#5)$) grid (#2);
  % }
  
  %   \foreach \i in {2,...,9} {
  %   \coordinate (P) at ($(0,0)!{log10(\i)}!(S)$) {};
  %   \draw[shift=(P), xstep=#3cm, ystep=#4cm, line width=0.6mm, #6!75] ($(#1)+(0,0.5*#5)$) grid (#2);
  % }
  %   \draw[xstep=#3cm, ystep=#4cm, line width=0.8mm, #6!100] (#1) grid (#2);
  %   \draw[step=1mm, line width=0.1mm, #6!25] (#1) grid (#2);
  %   \draw[step=5mm, line width=0.2mm, #6!50] (#1) grid (#2);
}



\newcommand{\millimetrique}[2]{
  \draw[step=1mm, line width=0.1mm, brown!25] (#1) grid (#2);
  \draw[step=5mm, line width=0.2mm, brown!50] (#1) grid (#2);
  \draw[step=1cm, line width=0.4mm, brown!75] (#1) grid (#2);
  \draw[step=5cm, line width=0.6mm, brown!100] (#1) grid (#2);
}





\newcommand{\cercle}[2]{
  \draw[fill=none, line width=0.6mm, brown!100](#1) circle (#2) ;
  
  \pgfmathsetmacro{\debut}{90}
  \pgfmathsetmacro{\pas}{360./2./3.141592654/#2)}
  \pgfmathsetmacro{\un}{\debut + \pas}
  \pgfmathsetmacro{\demi}{\debut + \pas/2.}
  \pgfmathsetmacro{\dixieme}{\debut + \pas/10.}
  \pgfmathsetmacro{\fin}{\debut + 360}
  
  % \foreach \i [count=\k from 0] in {0,11.4591559026,...,359} {
  % \foreach \i [count=\k from 0] in {90, 101.4591559026,...,449} {
  \foreach \i [count=\k from 0] in {\debut, \un,...,\fin} {
    \draw ($(#1)+(\i:#2+.6)$) node[brown] {\textbf{\k}};
    \draw[line width=0.4mm, brown!100] ($(#1)+(\i:#2-.4)$) -- ($(#1)+(\i:#2+.4)$);
  }
  
  \foreach \i [count=\k from 0] in {\debut, \demi,...,\fin} {
    \draw[line width=0.4mm, brown!100] ($(#1)+(\i:#2-.3)$) -- ($(#1)+(\i:#2+.3)$);
  }
  
  \foreach \i [count=\k from 0] in {\debut, \dixieme,...,\fin} {
    \draw[line width=0.2mm, brown!100] ($(#1)+(\i:#2-.15)$) -- ($(#1)+(\i:#2+.15)$);
  }
}
